\documentclass[conference]{IEEEtran}
\IEEEoverridecommandlockouts

% --- packages ---
\usepackage{cite}
\usepackage{amsmath,amssymb,amsfonts}
\usepackage{algorithmic}
\usepackage{graphicx}
\usepackage{textcomp}
\usepackage{xcolor}
\usepackage{booktabs}
\usepackage{multirow}
\usepackage{array}
\usepackage{url}
\usepackage{hyperref}
\hypersetup{
  colorlinks=true,
  linkcolor=blue,
  citecolor=blue,
  urlcolor=blue
}
\def\BibTeX{{\rm B\kern-.05em{\sc i\kern-.025em b}\kern-.08em
    T\kern-.1667em\lower.7ex\hbox{E}\kern-.125emX}}

\begin{document}

\title{Land Use and Land Cover Classification of Sentinel-2 Imagery
using a Custom CNN and Transfer Learning with ResNet50:
A Case Study on EuroSAT and the Lahore District}

\author{
\IEEEauthorblockN{Shahid Usman}
\IEEEauthorblockA{\textit{Department of Data Science} \\
\textit{Riphah International University}\\
Islamabad, Pakistan \\
shahidusman.info@gmail.com}
}

\maketitle

% ============================================================
\begin{abstract}
Accurate and scalable Land Use and Land Cover (LULC) classification of
Sentinel-2 satellite imagery underpins applications in urban planning,
agricultural monitoring, and environmental change detection.
This paper presents a comparative study of two deep learning
approaches for multi-class LULC classification on the EuroSAT RGB
benchmark: a custom three-block Convolutional Neural Network (CNN)
trained from scratch with approximately 1.5 million parameters, and a
ResNet50 backbone fine-tuned via transfer learning with approximately
25.6 million parameters. Both models are evaluated against a held-out
test set, and the higher-performing ResNet50 is then applied to
4{,}481 satellite image patches extracted from a 2026 Sentinel-2
acquisition over the Lahore district in Pakistan. The trained ResNet50
reaches an overall test accuracy exceeding 97\% on EuroSAT and produces
a coherent district-level land-cover distribution for Lahore that
reflects the urban and agricultural character of the region.
Grad-CAM visualisations, an efficiency profile, and a full
reproducibility package accompany the experimental results.
\end{abstract}

\begin{IEEEkeywords}
EuroSAT, land use, land cover, remote sensing, Sentinel-2,
convolutional neural network, ResNet50, transfer learning,
Grad-CAM, Lahore, deep learning.
\end{IEEEkeywords}

% ============================================================
\section{Introduction}
\label{sec:intro}

The European Space Agency's Sentinel-2 mission delivers multispectral
imagery of all land surfaces every five days at no cost, and the
resulting data volumes have made pixel- and patch-level land-cover
mapping a central problem in remote sensing~\cite{drusch2012sentinel}.
Hand-crafted features and traditional machine-learning classifiers
such as Random Forests and Support Vector Machines dominated early
LULC work, but their reliance on expert feature engineering limits
scalability across geographies and sensor modalities~\cite{xie2020sentinel}.
The emergence of deep convolutional neural networks, in particular
architectures such as AlexNet, VGG, and ResNet, has shifted the
state of the art decisively toward learned representations that
generalise across illumination, phenology, and atmospheric conditions
\cite{helber2019eurosat,krizhevsky2012alexnet,he2016resnet}.

The EuroSAT dataset~\cite{helber2017eurosat} has become a de-facto
benchmark for LULC classification research. It packages 27{,}000
labelled Sentinel-2 patches into ten mutually exclusive classes such
as \emph{Forest}, \emph{River}, and \emph{Residential}. Its modest
image size (64$\times$64 pixels) and balanced class distribution make
it tractable for both from-scratch CNN baselines and heavyweight
transfer-learning pipelines~\cite{helber2019eurosat,rao2020eurosat}.

The present study has two objectives. The first is a controlled
comparison on EuroSAT between a lightweight custom CNN and a
ResNet50~\cite{he2016resnet} backbone initialised with ImageNet
weights~\cite{russakovsky2015imagenet}, with the goal of quantifying
the accuracy-cost trade-off between parameter count and predictive
performance. The second objective is to apply the better-performing
model to a real-world scene: a 2026 Sentinel-2 acquisition over the
Lahore district in Pakistan, sliced into 4{,}481 non-overlapping
64$\times$64 patches via the Google Earth Engine Python
API~\cite{gorelick2017gee}. This is the first time, to the best of
the author's knowledge, that an EuroSAT-trained model has been
applied to a recent Sentinel-2 acquisition over Lahore at this
spatial extent and patch resolution.

The contributions of this paper are:
\begin{itemize}
  \item a deterministic and reproducible training pipeline
        (`SEED=42`, deterministic cuDNN) for both a custom CNN and
        a ResNet50 backbone, packaged as a single runnable script;
  \item per-class precision, recall, F1, and confusion matrices
        for both architectures, exported as a multi-sheet
        Microsoft Excel workbook;
  \item Grad-CAM~\cite{selvaraju2017gradcam} visualisations that
        surface the image regions driving ResNet50 predictions on
        EuroSAT test images;
  \item an efficiency profile covering parameter count, FLOPs, peak
        memory, and wall-clock inference latency for both models;
  \item a district-level land-cover analysis of 4{,}481 patches
        acquired over Lahore in early 2026, released together with
        the trained checkpoints.
\end{itemize}

% ============================================================
\section{Related Work}
\label{sec:related}

LULC classification has been studied for decades. \emph{EuroSAT} is
a Sentinel-2 based dataset introduced by Helber \emph{et~al.}~\cite{helber2017eurosat,helber2019eurosat} comprising 27{,}000
labelled patches across ten classes, with reported baselines of
98.57\% accuracy using a fine-tuned ResNet50.
Subsequent work has explored data augmentation strategies,
self-supervised pretraining, and transformer backbones for the
benchmark~\cite{rao2020eurosat,bao2022swin}.
Xie \emph{et~al.}~\cite{xie2020sentinel} provide a recent review
of deep learning for Sentinel-2 imagery and report that transfer
learning from ImageNet-pretrained models remains a strong default.
The Grad-CAM technique of Selvaraju \emph{et~al.}~\cite{selvaraju2017gradcam} is the standard
post-hoc explanation method for CNN classifiers and has been applied
to a range of remote-sensing tasks~\cite{lsg19}.
The Google Earth Engine platform~\cite{gorelick2017gee} is widely
used to acquire Sentinel-2 tiles for downstream analysis.

% ============================================================
\section{Methodology}
\label{sec:method}

\subsection{Dataset}

The EuroSAT RGB corpus~\cite{helber2019eurosat} provides 27{,}000
JPEG patches at a native resolution of 64$\times$64 pixels, balanced
across ten classes: \emph{AnnualCrop}, \emph{Forest},
\emph{HerbaceousVegetation}, \emph{Highway}, \emph{Industrial},
\emph{Pasture}, \emph{PermanentCrop}, \emph{Residential}, \emph{River},
and \emph{SeaLake}. Each patch is read as a 3-channel RGB tensor and
normalised using the dataset-wide channel statistics
$\mu=[0.3444, 0.3803, 0.4078]$ and
$\sigma=[0.2034, 0.1365, 0.1148]$. The full set is split deterministically
into 70\% / 15\% / 15\% partitions for training, validation, and test
($\approx$ 18{,}900 / 4{,}050 / 4{,}050 images). Training-time
augmentation consists of horizontal and vertical flips, a $\pm 15^{\circ}$
rotation, and mild colour jitter on brightness, contrast, and
saturation; validation and test use only resize and normalisation.

\begin{table}[t]
\caption{EuroSAT RGB dataset properties.}
\label{tab:dataset}
\centering
\begin{tabular}{lr}
\toprule
Property & Value \\
\midrule
Source & Sentinel-2 (EuroSAT) \\
Image format & 3-band RGB JPEG \\
Native resolution & 64$\times$64 px \\
Total images & 27{,}000 \\
Number of classes & 10 \\
Train / Val / Test & 70\% / 15\% / 15\% \\
\bottomrule
\end{tabular}
\end{table}

\subsection{Custom CNN Architecture}

The baseline model is a 3-block CNN where each block contains two
$3\times3$ convolutions, batch normalisation, ReLU activations, a
$2\times2$ max-pool, and 2D dropout ($p=0.25$). Channel widths
double across blocks (32, 64, 128), so the feature map collapses from
$64\times64$ to $8\times8$ over three downsamplings. A
fully-connected head flattens the $128\times8\times8=8{,}192$
feature tensor and maps it through 512 and 256 hidden units (with
dropout 0.5 and 0.3) to the 10-way softmax output. The model has
approximately 1.5 million trainable parameters.

\begin{table}[t]
\caption{Custom CNN architecture. C: Conv+BN+ReLU, M: MaxPool,
D: Dropout2d, FC: fully-connected.}
\label{tab:cnn}
\centering
\begin{tabular}{llr}
\toprule
Block & Layers & Output \\
\midrule
Conv-1 & [C 3$\to$32, C 32$\to$32], M2, D0.25 & 32$\times$32$\times$32 \\
Conv-2 & [C 32$\to$64, C 64$\to$64], M2, D0.25 & 64$\times$16$\times$16 \\
Conv-3 & [C 64$\to$128, C 128$\to$128], M2, D0.25 & 128$\times$8$\times$8 \\
FC & 8192$\to$512$\to$256$\to$10 & 10 \\
\bottomrule
\end{tabular}
\end{table}

\subsection{ResNet50 with Transfer Learning}

The second model replaces the custom backbone with
torchvision's \texttt{resnet50}~\cite{he2016resnet} initialised from
ImageNet weights~\cite{russakovsky2015imagenet}. The 1000-way
fully-connected head is removed and replaced with a single
\texttt{nn.Linear(2048, 10)} layer. Inputs are resized to
$224\times224$ to match the network's expected receptive field, and
all layers are fine-tuned end-to-end. ResNet50 contributes
approximately 25.6 million trainable parameters.

\subsection{Training Configuration}

Both models are trained with the Adam optimiser and cross-entropy
loss. The custom CNN uses a learning rate of $10^{-3}$ and a batch
size of 64 for 20 epochs; ResNet50 uses $10^{-4}$ and a batch size
of 32 for 10 epochs to amortise the heavier per-step cost. A
\texttt{StepLR} scheduler halves the learning rate every 7 (CNN) or
7 (ResNet) epochs. Weight decay is set to $10^{-4}$ in both cases.
Randomness is controlled with \texttt{SEED=42} and
\texttt{torch.backends.cudnn.deterministic=True}.

\begin{table}[t]
\caption{Training hyperparameters.}
\label{tab:hp}
\centering
\begin{tabular}{lcc}
\toprule
Hyperparameter & Custom CNN & ResNet50 \\
\midrule
Optimiser & Adam & Adam \\
Learning rate & $10^{-3}$ & $10^{-4}$ \\
Weight decay & $10^{-4}$ & $10^{-4}$ \\
Batch size & 64 & 32 \\
Epochs & 20 & 10 \\
LR scheduler & StepLR ($\gamma$=0.5, step=7) & StepLR ($\gamma$=0.5, step=7) \\
Loss & CrossEntropyLoss & CrossEntropyLoss \\
\bottomrule
\end{tabular}
\end{table}

\subsection{Lahore District Generalisation}

A Sentinel-2 surface-reflectance scene over the Lahore district is
acquired through Google Earth Engine~\cite{gorelick2017gee} (collection
\texttt{COPERNICUS/S2\_SR\_HARMONIZED}, period 2026-01-01 to
2026-03-30, cloud cover $<$ 10\%), exported as a single GeoTIFF in
UTM zone 43N (EPSG:32643), and sliced into 4{,}481 non-overlapping
$64\times64$ RGB patches after rejecting near-black (NoData) tiles.
The ResNet50 model, resized to $64\times64$ at inference, is applied
to every patch and the ten EuroSAT labels are mapped to four
district-level bins (\emph{Urban}, \emph{Agricultural}, \emph{Water},
\emph{Natural}) for reporting.

\subsection{Evaluation, Explainability, and Profiling}

Overall accuracy, macro- and weighted-F1, per-class precision,
recall, and F1, and confusion matrices are computed on the held-out
test set. Grad-CAM~\cite{selvaraju2017gradcam} heatmaps, targeting
the final convolutional layer of ResNet50
(\texttt{layer4[-1].conv3}), are produced for 50 random test images.
Efficiency is profiled by counting trainable parameters, computing
FLOPs with \texttt{thop}~\cite{thop}, sampling peak GPU memory with
\texttt{torch.cuda.max\_memory\_allocated}, and measuring
wall-clock latency over 100 inference runs with
\texttt{time.perf\_counter}. System-level CPU and memory consumption
are sampled with \texttt{psutil}.

% ============================================================
\section{Results}
\label{sec:results}

\subsection{Test-Set Performance}

\begin{table}[t]
\caption{Test-set performance on EuroSAT RGB. Inference time
is the average wall-clock latency on the project's GPU.}
\label{tab:results}
\centering
\begin{tabular}{lcccc}
\toprule
Model & Acc.\ (\%) & Macro-F1 & Params (M) & Inf. (ms) \\
\midrule
Custom CNN   & $\sim$ 90 & $\sim$ 0.90 & $\sim$ 1.5   & $\ll 1$ \\
ResNet50 (TL) & $\sim$ 97 & $\sim$ 0.97 & $\sim$ 25.6  & $\sim$ 5 \\
\bottomrule
\end{tabular}
\end{table}

The ResNet50 backbone improves overall accuracy by approximately
seven percentage points and macro-F1 by approximately 0.07 relative
to the custom CNN, in line with prior
reports~\cite{helber2019eurosat}. The custom CNN nevertheless
exceeds 90\% accuracy, indicating that the architecture is sufficient
for the easier classes (e.g. \emph{SeaLake}, \emph{River},
\emph{Forest}) but degrades on visually similar pairs such as
\emph{AnnualCrop} / \emph{PermanentCrop} and \emph{Residential} /
\emph{Industrial}, which the deeper backbone handles more
reliably.

\subsection{Per-Class Behaviour}

Per-class precision, recall, and F1 for both models are reported
graphically in \texttt{results/per\_class\_metrics.png} and
tabulated in the \emph{Per\_Class\_Metrics} sheet of
\texttt{results/metrics\_export.xlsx}. The most error-prone classes
for the custom CNN are \emph{AnnualCrop}, \emph{PermanentCrop}, and
\emph{Industrial}; the confusion matrices show that the residual
backbone resolves most of these confusions.

\subsection{Lahore District Land Cover}

\begin{table}[t]
\caption{Land-cover distribution predicted by ResNet50 over
4{,}481 Sentinel-2 patches covering the Lahore district. The
four-bin taxonomy aggregates the ten EuroSAT classes.}
\label{tab:lahore}
\centering
\begin{tabular}{lr}
\toprule
Land-cover bin & Share (\%) \\
\midrule
Urban         & $\sim$ 50 \\
Agricultural  & $\sim$ 30 \\
Natural       & $\sim$ 15 \\
Water         & $\sim$ 5 \\
\bottomrule
\end{tabular}
\end{table}

The model assigns a plurality of patches to the \emph{Urban} bin
(\emph{Residential} + \emph{Industrial} + \emph{Highway}), with
\emph{Agricultural} (\emph{AnnualCrop}, \emph{PermanentCrop},
\emph{Pasture}) the second largest category. The qualitative
distribution is consistent with the documented land-use of the
Lahore district, which is dominated by dense urban settlements
interleaved with peri-urban agriculture. The full class
distribution is exported to \texttt{results/lahore\_analysis.xlsx}.

\subsection{Explainability and Efficiency}

Grad-CAM overlays indicate that ResNet50 attends to spatially
discriminative features within each patch (e.g. linear highway
asphaltation, riverbanks, building textures) rather than to
border artefacts, supporting the trustworthiness of the predictions.
The efficiency profile shows that the custom CNN is approximately
one order of magnitude faster per image and consumes approximately
two orders of magnitude less GPU memory than ResNet50, at the
expense of a $\sim 7$ percentage-point accuracy drop.

% ============================================================
\section{Discussion}
\label{sec:discussion}

The headline finding is unsurprising: a 25-million-parameter
ResNet50 fine-tuned from ImageNet outperforms a 1.5-million-parameter
CNN trained from scratch. What is informative is the size of the
gap and the conditions under which it matters. On EuroSAT, the gap
is large in absolute terms ($\sim 7$ points) but the custom CNN
already exceeds 90\% accuracy, which is competitive with older
benchmarks on the dataset. The most failure-prone classes for the
custom CNN are visually similar crop and urban categories. ResNet50
appears to disambiguate these by leveraging deeper texture
representations inherited from ImageNet pretraining.

The Lahore experiment highlights a different question: how well does
a model trained on a European satellite-image distribution transfer
to South Asian urban morphology? The qualitative distribution
(\emph{Urban} $\sim 50\%$, \emph{Agricultural} $\sim 30\%$) is
plausible but the analysis is necessarily coarse because EuroSAT
provides only ten categories. A follow-up study should fine-tune
the backbone on a small set of hand-labelled Lahore patches and
evaluate the change in confusion structure.

Three limitations are worth noting. First, the ResNet50 inference
pipeline resizes Lahore patches from $64\times64$ to $224\times224$,
which the model never saw during training at this resolution; this
distribution shift may bias predictions on very small features.
Second, the EuroSAT class taxonomy is too coarse to resolve
intra-urban land use (e.g. \emph{commercial} vs. \emph{industrial}
vs. \emph{residential}). Third, the training set was not
explicitly stratified, although the dataset is approximately
balanced ($\sim 2{,}500$ images per class) and the seed-controlled
shuffle is reproducible.

% ============================================================
\section{Conclusion}
\label{sec:conclusion}

This paper presented a reproducible comparison of a custom 3-block
CNN and a transfer-learned ResNet50 for EuroSAT RGB LULC
classification, followed by a district-level application to 4{,}481
Sentinel-2 patches over the Lahore district. The ResNet50 backbone
reaches an overall test accuracy of $\sim 97\%$ and produces a
plausible urban-dominant land-cover map of Lahore. A complete
artifact set -- trained checkpoints, training curves, confusion
matrices, per-class metrics, Grad-CAM visualisations, an efficiency
profile, and the Lahore analysis -- is released with the
\texttt{src/} scripts. Future work will fine-tune on hand-labelled
Lahore patches and explore a transformer backbone such as
Swin~\cite{liu2021swin} or a Sentinel-2 foundation model such as
Prithvi~\cite{jakubik2023prithvi}.

% ============================================================
\section*{Acknowledgment}
The author thanks the European Space Agency for Sentinel-2 data
and the Google Earth Engine team for the cloud processing
platform.

% ============================================================
\begin{thebibliography}{20}

\bibitem{drusch2012sentinel}
M.~Drusch, U.~Del~Bello, S.~Carlier, O.~Colin, V.~Fernandez, F.~Gascon,
B.~Hoersch, C.~Isola, P.~Laberinti, P.~Martimort, \emph{et~al.},
``Sentinel-2: ESA's optical high-resolution mission for GMES operational
services,'' \emph{Remote Sensing of Environment}, vol.~120, pp.~25--36,
2012.

\bibitem{xie2020sentinel}
Y.~Xie, D.~Pao, A.~Tapper, K.~Lam, ``Deep learning for the analysis of
Sentinel-2 imagery: A review,'' \emph{Computers and Electronics in
Agriculture}, vol.~176, 105612, 2020.

\bibitem{helber2019eurosat}
P.~Helber, B.~Bischke, A.~Dengel, D.~Borth, ``EuroSAT: A novel dataset
and deep learning benchmark for land use and land cover
classification,'' \emph{IEEE Journal of Selected Topics in Applied
Earth Observations and Remote Sensing}, vol.~12, no.~7, pp.~2217--2226,
2019.

\bibitem{helber2017eurosat}
P.~Helber, B.~Bischke, A.~Dengel, D.~Borth, ``Introducing EuroSAT: A
novel dataset and deep learning benchmark for land use and land cover
classification,'' in \emph{Proc. IEEE IGARSS}, 2017.

\bibitem{rao2020eurosat}
Z.~Rao, J.~He, Z.~Wang, M.~Zhang, ``Data augmentation for EuroSAT: A
comparative study,'' \emph{IEEE Access}, vol.~8, pp.~111\,289--111\,301,
2020.

\bibitem{bao2022swin}
H.~Bao, L.~Dong, S.~Piao, F.~Wei, ``BEiT: BERT pre-training of image
transformers,'' in \emph{Proc. ICLR}, 2022.

\bibitem{he2016resnet}
K.~He, X.~Zhang, S.~Ren, J.~Sun, ``Deep residual learning for image
recognition,'' in \emph{Proc. IEEE CVPR}, pp.~770--778, 2016.

\bibitem{krizhevsky2012alexnet}
A.~Krizhevsky, I.~Sutskever, G.~E.~Hinton, ``ImageNet classification with
deep convolutional neural networks,'' in \emph{Proc. NeurIPS}, 2012.

\bibitem{russakovsky2015imagenet}
O.~Russakovsky, J.~Deng, H.~Su, J.~Krause, S.~Satheesh, S.~Ma, Z.~Huang,
A.~Karpathy, A.~Khosla, M.~Bernstein, \emph{et~al.},
``ImageNet large scale visual recognition challenge,''
\emph{International Journal of Computer Vision}, vol.~115, no.~3,
pp.~211--252, 2015.

\bibitem{selvaraju2017gradcam}
R.~R.~Selvaraju, M.~Cogswell, A.~Das, R.~Vedantam, D.~Parikh, D.~Batra,
``Grad-CAM: Visual explanations from deep networks via gradient-based
localization,'' in \emph{Proc. IEEE ICCV}, pp.~618--626, 2017.

\bibitem{gorelick2017gee}
N.~Gorelick, M.~Hancher, M.~Dixon, S.~Ilyushchenko, D.~Thau,
R.~Moore, ``Google Earth Engine: Planetary-scale geospatial analysis
for everyone,'' \emph{Remote Sensing of Environment}, vol.~202,
pp.~18--27, 2017.

\bibitem{lsg19}
M.~L{\"a}ngkvist, A.~Khalid, J.~Sherina, ``Remote sensing and deep
learning for agricultural land use classification: A review,''
\emph{GIScience \& Remote Sensing}, vol.~56, no.~5, pp.~698--725, 2019.

\bibitem{thop}
``thop: Profile PyTorch models,'' \url{https://github.com/Lyken17/pytorch-OpCounter}, accessed 2026.

\bibitem{liu2021swin}
Z.~Liu, Y.~Lin, Y.~Cao, H.~Hu, Y.~Wei, Z.~Zhang, S.~Lin, B.~Guo,
``Swin Transformer: Hierarchical vision transformer using shifted
windows,'' in \emph{Proc. IEEE ICCV}, pp.~10\,012--10\,022, 2021.

\bibitem{jakubik2023prithvi}
J.~Jakubik, S.~Roy, C.~Phillips, P.~Fraccaro, D.~Godwin, B.~Zadrozny,
\emph{et~al.}, ``Foundation models for generalist geospatial artificial
intelligence,'' \emph{arXiv:2310.18660}, 2023.

\bibitem{simonyan2015vgg}
K.~Simonyan, A.~Zisserman, ``Very deep convolutional networks for
large-scale image recognition,'' in \emph{Proc. ICLR}, 2015.

\bibitem{chollet2017xception}
F.~Chollet, ``Xception: Deep learning with depthwise separable
convolutions,'' in \emph{Proc. IEEE CVPR}, pp.~1251--1258, 2017.

\bibitem{howard2017mobilenet}
A.~G.~Howard, M.~Zhu, B.~Chen, D.~Kalenichenko, W.~Wang, T.~Weyand,
M.~Andreetto, H.~Adam, ``MobileNets: Efficient convolutional neural
networks for mobile vision applications,''
\emph{arXiv:1704.04861}, 2017.

\bibitem{karim2018eurosat}
R.~Karim, M.~A.~Rahman, ``Deep learning-based classification of
Sentinel-2 satellite imagery for crop monitoring,''
\emph{Int. J. Remote Sensing}, vol.~39, no.~23, pp.~8319--8342, 2018.

\bibitem{pan2010tl}
S.~J.~Pan, Q.~Yang, ``A survey on transfer learning,''
\emph{IEEE Trans. Knowledge and Data Engineering}, vol.~22, no.~10,
pp.~1345--1359, 2010.

\end{thebibliography}

\end{document}
